\documentclass[12pt]{article}
\usepackage{amsmath}
\usepackage{amssymb}
\usepackage{amsthm}


\newtheorem{theorem}{Theorem}

\begin{document}
\title{Submission D solution to Problem 3}
\date{}
\maketitle


\section*{Problem Statement}

Let $w_1, \ldots, w_m$ be non-negative real numbers such that $\sum_{i=1}^m w_i=1$, and let $v_1, \ldots, v_m$ be i.i.d.\ Bernoulli$(p)$ random variables. For which values of probability $p$ is it true that
\[\Pr\Big[\sum_{i=1}^m w_i v_i\geq p\Big]\geq p?\]

\section*{Partial Solution}

\textit{Note: The following represents a partial solution to the problem. As noted at the end of the proof, sufficiency is fully established, but there remains a structural gap in proving the necessity for non-reciprocal values of $p$.}

\begin{theorem}
Let $w_1, \ldots, w_m$ be non-negative real numbers such that $\sum_{i=1}^m w_i = 1$, and let $v_1, \ldots, v_m$ be i.i.d.\ Bernoulli$(p)$ random variables. The inequality 
\[\Pr\Big[\sum_{i=1}^m w_i v_i\geq p\Big]\geq p\]
holds for all valid weight distributions if and only if $p = 1/k$ for some integer $k \ge 1$ (or $p=0$).
\end{theorem}

\begin{proof}
For $p = 0$, the condition reduces to $\Pr[0 \ge 0] \ge 0$, which evaluates to $1 \ge 0$ and is trivially true.

Assume $p = 1/k$ for some integer $k \ge 1$. 
We construct $k$ coupled sequences of random variables $v_i^{(j)}$ for indices $i \in \{1, \dots, m\}$ and $j \in \{1, \dots, k\}$ as follows:
For each $i$, select an index $J_i$ uniformly at random from the set $\{1, \dots, k\}$, such that the choices $J_1, \ldots, J_m$ are mutually independent. 
Set $v_i^{(J_i)} = 1$, and set $v_i^{(j)} = 0$ for all $j \neq J_i$.

By this construction, for any fixed sequence index $j$, the variables $v_1^{(j)}, \ldots, v_m^{(j)}$ are independent Bernoulli$(1/k)$ random variables. 
Define $X^{(j)} = \sum_{i=1}^m w_i v_i^{(j)}$. 
For each $j$, the random variable $X^{(j)}$ has the exact same distribution as the original sum $\sum_{i=1}^m w_i v_i$.

Summing across all $k$ sequences yields:
\[\sum_{j=1}^k X^{(j)} = \sum_{j=1}^k \sum_{i=1}^m w_i v_i^{(j)} = \sum_{i=1}^m w_i \sum_{j=1}^k v_i^{(j)}\]
Since exactly one $v_i^{(j)} = 1$ for each $i$ (specifically when $j = J_i$), the inner sum $\sum_{j=1}^k v_i^{(j)}$ equals $1$ deterministically.
Thus,
\[\sum_{j=1}^k X^{(j)} = \sum_{i=1}^m w_i (1) = 1\]

Because the $k$ identically distributed random variables $X^{(j)}$ sum to exactly $1$, at least one of them must evaluate to $1/k$ or greater. 
Let $A_j$ be the event that $X^{(j)} \ge 1/k$. We have $\Pr\Big[\bigcup_{j=1}^k A_j\Big] = 1$.
By the union bound:
\[ 1 = \Pr\Big[\bigcup_{j=1}^k A_j\Big] \le \sum_{j=1}^k \Pr[A_j] = k \Pr\Big[X^{(1)} \ge \frac{1}{k}\Big] \]
Dividing by $k$ gives:
\[ \Pr\Big[\sum_{i=1}^m w_i v_i \ge \frac{1}{k}\Big] \ge \frac{1}{k} \]
This establishes sufficiency.

For values of $p \notin \{0\} \cup \{1/k \mid k \in \mathbb{Z}^+\}$, the inequality does not universally hold.

\vspace{\baselineskip}
\noindent\textbf{[Gap: The explicit construction of pathological weight distributions providing counterexamples for all non-reciprocal values of $p$ is structurally omitted.]}
\end{proof}

\end{document}